\documentclass[11pt,a4paper]{article}
\usepackage[margin=19mm,top=17mm,bottom=17mm]{geometry}
\usepackage{amsmath,amssymb,amsfonts}
\usepackage{mathptmx}
\usepackage{enumitem}
\usepackage{fancyhdr}
\usepackage{draftwatermark}
\usepackage{graphicx}
\usepackage{array,tabularx,booktabs}
\usepackage[hidelinks]{hyperref}
\setlength{\parindent}{0pt}\setlength{\parskip}{4pt}
\setlist[enumerate,1]{leftmargin=1.1em,itemsep=3pt,parsep=2pt}
\setlist[itemize]{leftmargin=1.4em,itemsep=1pt}
\pagestyle{fancy}\fancyhf{}
\fancyhead[L]{\footnotesize \textbf{}}
\fancyhead[C]{\footnotesize \textbf{Practice Paper}}
\fancyhead[R]{\footnotesize \textbf{}}
\fancyfoot[L]{\footnotesize Think C}
\fancyfoot[C]{\footnotesize Education is not just about Information  $\cdot$  Page \thepage}
\fancyfoot[R]{\footnotesize }
\renewcommand{\headrulewidth}{0.4pt}
\SetWatermarkText{ThinkC}
\SetWatermarkScale{1.4}
\SetWatermarkAngle{45}
\SetWatermarkLightness{82}
\begin{document}
\vspace{2mm}
\noindent \textbf{Marks : 100} \hfill \textbf{Time : 3 : 00 hours}
\noindent\rule{\linewidth}{1.1pt}\vspace{-1mm}\noindent\rule{\linewidth}{0.5pt}
{\centering \Large \textbf{STRAIGHT LINES 25 QUESTION SET}\par}
\noindent\rule{\linewidth}{1.1pt}\vspace{-1mm}\noindent\rule{\linewidth}{0.5pt}
\vspace{2mm}

\begin{enumerate}[resume]
\item In the Cartesian plane, the vertices of a triangle are $A(0,0)$, $B(6,0)$ and $C(2,4)$. The orthocentre of $	riangle ABC$ is \hfill [1 marks]
{\renewcommand{\arraystretch}{1.8}
\begin{tabularx}{\linewidth}{X X X X}
(A) \quad $(2,2)$ & (B) \quad $(2,4)$ & (C) \quad $(4,2)$ & (D) \quad $(0,2)$ \\[10pt]
\end{tabularx}}
\item A straight line through the origin passes through the common point of the lines $3x+2y-13=0$ and $x-y-1=0$. Its slope is \hfill [1 marks]
{\renewcommand{\arraystretch}{1.8}
\begin{tabularx}{\linewidth}{X X X X}
(A) \quad $\frac{2}{3}$ & (B) \quad $\frac{3}{2}$ & (C) \quad $-\frac{2}{3}$ & (D) \quad $-\frac{3}{2}$ \\[10pt]
\end{tabularx}}
\item For each real value of $t$, consider the line
$$ (3+2t)x+(3-2t)y=12. $$
All these lines pass through a fixed point $P$. If the distance from $(u,v)$ to the origin is defined by
$$\rho((u,v),O)=\max\{|u|,|v|\},$$
then $\rho(P,O)$ is equal to: \hfill [1 marks]
{\renewcommand{\arraystretch}{1.8}
\begin{tabularx}{\linewidth}{X X X X}
(A) \quad $1$ & (B) \quad $2$ & (C) \quad $\sqrt{8}$ & (D) \quad $0$ \\[10pt]
\end{tabularx}}
\item Let the coefficients $a$, $b$ and $c$ of the variable line $ax+by+c=0$ be respectively the $2^{\text{nd}}$, $5^{\text{th}}$ and $9^{\text{th}}$ terms of an increasing arithmetic progression. This line passes through a fixed point. The fixed point lies on the curve \hfill [4 marks]
{\renewcommand{\arraystretch}{1.8}
\begin{tabularx}{\linewidth}{X X}
(A) \quad $x^2+y^2=\dfrac{65}{9}$ & (B) \quad $3x+4y=5$ \\[10pt]
(C) \quad $x^2+y^2=\dfrac{25}{9}$ & (D) \quad $4x-3y=9$ \\[10pt]
\end{tabularx}}
\item The area of the triangle enclosed by the line $2x+3y=12$ and the two angle bisectors of the pair of straight lines $(x+y-6)(x-y-2)=0$ is: \hfill [4 marks]
{\renewcommand{\arraystretch}{1.8}
\begin{tabularx}{\linewidth}{X X}
(A) \quad $\dfrac{1}{3}\text{ sq. units}$ & (B) \quad $\dfrac{2}{3}\text{ sq. units}$ \\[10pt]
(C) \quad $1\text{ sq. unit}$ & (D) \quad $\dfrac{4}{3}\text{ sq. units}$ \\[10pt]
\end{tabularx}}
\item Let $m$ be an integer. The number of values of $m$ for which the graph of $y=|mx-3|$ has at least one point of intersection with the line $5x+2y-11=0$ whose $x$-coordinate is an integer, is \hfill [4 marks]
{\renewcommand{\arraystretch}{1.8}
\begin{tabularx}{\linewidth}{X X X X}
(A) \quad $4$ & (B) \quad $5$ & (C) \quad $6$ & (D) \quad $7$ \\[10pt]
\end{tabularx}}
\item A variable line through the point $P(0,-3)$ meets the lines $y=x-2$ and $y=-x-2$ at $A$ and $B$, respectively. If the segment $AB$ subtends a right angle at the origin, then the possible slope(s) of the variable line are: \hfill [1 marks]
{\renewcommand{\arraystretch}{1.8}
\begin{tabularx}{\linewidth}{X X X X}
(A) \quad $m=\pm\sqrt{2}$ & (B) \quad $m=\pm 2$ & (C) \quad $m=\pm\frac{1}{\sqrt{2}}$ & (D) \quad $m=2$ only \\[10pt]
\end{tabularx}}
\item A line through the fixed point $P(3,-2)$ cuts the positive $x$- and $y$-axes at $M$ and $N$, respectively. If $G(u,v)$ denotes the centroid of the triangle formed by $O(0,0)$, $M$, and $N$, then the locus of $G$ is given by: \hfill [4 marks]
{\renewcommand{\arraystretch}{1.8}
\begin{tabularx}{\linewidth}{X X X X}
(A) \quad $3uv-3v+2u=0$ & (B) \quad $3uv+3v-2u=0$ & (C) \quad $uv-3v+2u=0$ & (D) \quad $3uv-2v+3u=0$ \\[10pt]
\end{tabularx}}
\item Let the points $P(2,1)$ and $Q(0,3)$ be two vertices of a non-degenerate triangle $PQR$. If its orthocentre coincides with the origin, then the coordinates of $R$ are \hfill [4 marks]
{\renewcommand{\arraystretch}{1.8}
\begin{tabularx}{\linewidth}{X X X X}
(A) \quad $(1,1)$ & (B) \quad $(1,-1)$ & (C) \quad $(-1,1)$ & (D) \quad $(2,2)$ \\[10pt]
\end{tabularx}}
\item The line $y=2x-8$ is rotated about its point of intersection with the $x$-axis until it becomes perpendicular to the line $3x+4y-7=0$. The equation of the resulting line is \hfill [1 marks]
{\renewcommand{\arraystretch}{1.8}
\begin{tabularx}{\linewidth}{X X X X}
(A) \quad $4x-3y-16=0$ & (B) \quad $3x+4y-12=0$ & (C) \quad $4x+3y-16=0$ & (D) \quad $3x-4y-12=0$ \\[10pt]
\end{tabularx}}
\item A right triangle $ABC$ has $ngle C=90^irc$, and its two legs are parallel to the coordinate axes. The median from $A$ lies on the line $y=-2x+3$, while the median from $B$ lies on a line of the form $y=mx+6$. The value of $m$ is \hfill [1 marks]
{\renewcommand{\arraystretch}{1.8}
\begin{tabularx}{\linewidth}{X X X X}
(A) \quad $-8$ & (B) \quad $-\frac{1}{2}$ & (C) \quad $\frac{1}{2}$ & (D) \quad $8$ \\[10pt]
\end{tabularx}}
\item Let $A(1,1)$, $B(7,1)$ and $C(1,1+h)$, where $h>0$. If the medians drawn from $A$ and $B$ of $\triangle ABC$ are perpendicular, then the value of $h$ is: \hfill [4 marks]
{\renewcommand{\arraystretch}{1.8}
\begin{tabularx}{\linewidth}{X X X X}
(A) \quad $3\sqrt{2}$ & (B) \quad $6\sqrt{2}$ & (C) \quad $12\sqrt{2}$ & (D) \quad $18\sqrt{2}$ \\[10pt]
\end{tabularx}}
\item The pair of straight lines passing through the point $(1,2)$ and making an angle of $30^\circ$ with the line $y=x$ is represented by \hfill [4 marks]
{\renewcommand{\arraystretch}{1.8}
\begin{tabularx}{\linewidth}{X X}
(A) \quad $(y-2)^2-4(x-1)(y-2)+(x-1)^2=0$ & (B) \quad $(y-2)^2+4(x-1)(y-2)+(x-1)^2=0$ \\[10pt]
(C) \quad $(y-2)^2-2(x-1)(y-2)-(x-1)^2=0$ & (D) \quad $(y-2)^2-4(x-1)(y-2)-(x-1)^2=0$ \\[10pt]
\end{tabularx}}
\item The vertices of a triangle are $A(-2,1)$, $B(4,1)$ and $C(1,5)$. The coordinates of its orthocentre are: \hfill [1 marks]
{\renewcommand{\arraystretch}{1.8}
\begin{tabularx}{\linewidth}{X X}
(A) \quad $(1,\frac{13}{4})$ & (B) \quad $(1,\frac{7}{4})$ \\[10pt]
(C) \quad $(\frac{1}{2},\frac{13}{4})$ & (D) \quad $(1,4)$ \\[10pt]
\end{tabularx}}
\item A line passing through the origin intersects each of the two lines $3x-2y+5=0$ and $x+4y-7=0$ at one common point. The slope of this line is: \hfill [4 marks]
{\renewcommand{\arraystretch}{1.8}
\begin{tabularx}{\linewidth}{X X X X}
(A) \quad $-\dfrac{13}{3}$ & (B) \quad $\dfrac{13}{3}$ & (C) \quad $-\dfrac{3}{13}$ & (D) \quad $\dfrac{3}{13}$ \\[10pt]
\end{tabularx}}
\item For each real value of $t$, consider the straight line $(3+t)x+(2-t)y=10$. All these lines pass through a fixed point $P$. If the distance of a point $(x,y)$ from the origin is defined by $D(x,y)=\max\{|x|,|y|\}$, then $D(P)$ is \hfill [1 marks]
{\renewcommand{\arraystretch}{1.8}
\begin{tabularx}{\linewidth}{X X X X}
(A) \quad $1$ & (B) \quad $2$ & (C) \quad $\sqrt{5}$ & (D) \quad $5$ \\[10pt]
\end{tabularx}}
\item The area of the triangle enclosed by the line $2x+3y=12$ and the two angle bisectors of the pair of straight lines
\[
(x-y+1)(x+y-3)=0
\]
is: \hfill [4 marks]
{\renewcommand{\arraystretch}{1.8}
\begin{tabularx}{\linewidth}{X X}
(A) \quad $\dfrac{4}{3}$ sq. units & (B) \quad $\dfrac{8}{3}$ sq. units \\[10pt]
(C) \quad $\dfrac{10}{3}$ sq. units & (D) \quad $4$ sq. units \\[10pt]
\end{tabularx}}
\item For how many integral values of $m$ does the line $3x+2y-13=0$ have at least one common point with the graph of $y=|mx+1|$ whose $x$-coordinate is an integer? \hfill [4 marks]
{\renewcommand{\arraystretch}{1.8}
\begin{tabularx}{\linewidth}{X X X X}
(A) \quad $6$ & (B) \quad $7$ & (C) \quad $8$ & (D) \quad $9$ \\[10pt]
\end{tabularx}}
\end{enumerate}
\end{document}